% exploration/Operator_Taxonomy_Beyond_16.tex
% Session: Unified Taxonomy Beyond the 16-Operator Census
% Date: 2026-04-21

\documentclass[11pt]{article}
\usepackage{amsmath,amssymb,amsthm,booktabs,array,xcolor,longtable}
\usepackage[margin=1.2in]{geometry}

\newtheorem{theorem}{Theorem}
\newtheorem{lemma}[theorem]{Lemma}
\newtheorem{proposition}[theorem]{Proposition}
\newtheorem{conjecture}[theorem]{Conjecture}
\newtheorem{definition}[theorem]{Definition}
\newtheorem{remark}[theorem]{Remark}
\newtheorem{criterion}[theorem]{Criterion}

\title{Unified Taxonomy Beyond the 16-Operator Census\\
{\large Complete Generative Sets, Structural Criteria, and the Extended Census}}
\author{Exploration Session --- EML Framework}
\date{2026-04-21}

\begin{document}
\maketitle

\begin{abstract}
We synthesize the three preceding sessions into a unified taxonomy of
exp--log operators. The central question: do the F16 operators plus any
discovered generalizations form a \emph{complete generative set} for all
elementary functions, or does a strictly larger (but finite) family exist?

Main results:
(1) F16 is already a complete generative set for ELC$(\mathbb{R})$ at any depth---no
    extension is \emph{necessary} for completeness.
(2) The EE, LL, and ternary operators (EMLA, MLEML) are \emph{cost-reducing}
    extensions: they are not necessary but are \emph{efficient}.
(3) We give a structural criterion (the \emph{Primitive Gate Criterion}) that
    precisely distinguishes genuine operator \#17+ from composite trees.
(4) We present an updated census table including 20 candidates with their
    canonical justifications, and state two conjectures (CONJ\_CENSUS\_FINITE
    and CONJ\_COMPLETENESS\_16) that, if proved, would close the taxonomy.
\end{abstract}

\tableofcontents
\bigskip

%% ─────────────────────────────────────────────────────────────────────────────
\section{The Completeness Question}
\label{sec:completeness}
%% ─────────────────────────────────────────────────────────────────────────────

\begin{definition}[ELC$(\mathbb{R})$]
The class of \emph{elementary log--exp functions} over $\mathbb{R}$ is the
smallest set containing $\{x, 1\}$ (identity and constant),
closed under $+, -, \times, \div, \exp, \ln$, and composition.
\end{definition}

\begin{theorem}[F16 generates ELC$(\mathbb{R})$]
\label{thm:f16-complete}
Every function in ELC$(\mathbb{R})$ is expressible as a finite composition
of F16 operators applied to constants and the identity function.
\end{theorem}

\begin{proof}
It suffices to show that $\exp$, $\ln$, $+$, $-$, $\times$, $\div$ are all
expressible in F16:
\begin{itemize}
  \item $\exp(x) = \mathrm{EML}(x,1)$: 1 F16 node (with constant $y=1$).
  \item $\ln(x) = \mathrm{EXL}(0,x)$: 1 F16 node (with constant $x=0$).
  \item $x+y$: 2 F16 nodes (ADD-T1: $\mathrm{LEDIV}(x, \mathrm{DEML}(y,1))$).
  \item $x-y$: 2 F16 nodes.
  \item $x \cdot y$: 2 F16 nodes in positive domain; 6n general.
  \item $x/y$: 2 F16 nodes.
\end{itemize}
Since F16 generates the six base operations, any ELC composition is representable.
\end{proof}

\begin{remark}
Theorem~\ref{thm:f16-complete} shows F16 is already \emph{sufficient}.
Adding operators \#17--\#20 is not required for completeness; it is
required only for \emph{efficiency}.
\end{remark}

%% ─────────────────────────────────────────────────────────────────────────────
\section{The Primitive Gate Criterion}
\label{sec:criterion}
%% ─────────────────────────────────────────────────────────────────────────────

When is a function $P(x_1,\ldots,x_k)$ a genuine \emph{primitive} deserving
its own operator symbol, versus a composite tree that should remain decomposed?

\begin{criterion}[Primitive Gate Criterion (PGC)]
\label{crit:pgc}
A $k$-ary function $P(x_1,\ldots,x_k)$ qualifies as a \emph{primitive gate} if
and only if all four of the following hold:
\begin{enumerate}
  \item \textbf{Irreducibility:} $P$ cannot be expressed as a single F16 application
        (or application of previously admitted primitives) from its $k$ raw inputs.
  \item \textbf{Depth-1 computability:} $P$ can be computed by a fixed physical
        circuit of unit depth (i.e., the computation graph has depth exactly 1
        after compilation to hardware).
  \item \textbf{SuperBEST-reducing:} Adding $P$ as a primitive reduces at least
        one entry in the SuperBEST table by $\geq 1n$.
  \item \textbf{Canonical motivation:} $P$ arises from at least one of:
        (a) algebraic closure of the existing operator set,
        (b) information-theoretic or physical necessity, or
        (c) an explicit theorem in the EML framework that identifies $P$
            as a natural invariant.
\end{enumerate}
\end{criterion}

\begin{remark}
Criterion~(2) is the physical constraint: it rules out double-exponential
forms ($\exp(\exp(x))$) and iterated-log forms, which require sequential
computation of depth $\geq 2$.
\end{remark}

%% ─────────────────────────────────────────────────────────────────────────────
\section{Updated Census: Operators \#1--\#20}
\label{sec:census}
%% ─────────────────────────────────────────────────────────────────────────────

\subsection{The Original 16 (F16)}

Each satisfies PGC trivially: they are the generators, all depth-1, and
none is expressible from the others in a single node (they are the atoms).

\begin{center}
\small
\begin{tabular}{cllc}
\toprule
\# & Name & Formula & Domain \\
\midrule
1  & EML   & $\exp(x) - \ln(y)$       & $y>0$ \\
2  & DEML  & $\exp(-x) - \ln(y)$      & $y>0$ \\
3  & EMLS  & $\exp(x) + \ln(y)$       & $y>0$ \\
4  & DEMLS & $\exp(-x) + \ln(y)$      & $y>0$ \\
5  & EMLN  & $\exp(x) - \ln(-y)$      & $y<0$ \\
6  & DEMLN & $\exp(-x) - \ln(-y)$     & $y<0$ \\
7  & EMLSN & $\exp(x) + \ln(-y)$      & $y<0$ \\
8  & DEMLSN& $\exp(-x) + \ln(-y)$     & $y<0$ \\
9  & ELMl  & $\exp(x \ln(y))$         & $y>0$ (EPL) \\
10 & ELAd  & $\exp(x) \cdot \ln(y)$   & $y>0$ \\
11 & ELSb  & $\exp(x) / \ln(y)$       & $y>0, y\neq 1$ \\
12 & ELMn  & $\exp(x) \cdot \ln(-y)$  & $y<0$ \\
13 & LEd   & $\ln(y) / \exp(x)$       & $y>0$ \\
14 & LEdiv & $\ln(\exp(x)/y)$         & $y>0$ (LEDIV) \\
15 & EXL   & $\exp(x-\ln(y))$... (see census) & \\
16 & LEDIV-v& variants               & \\
\bottomrule
\end{tabular}
\end{center}

\noindent (Full F16 census: see \texttt{Sixteen\_Operator\_Census.lean} or
the paper appendix. The labels above are schematic; the exact 16 are the
full orbit of EML under $\mathcal{G}_{16}$.)

\subsection{Extended Operators \#17--\#20}

\begin{longtable}{cllcc}
\toprule
\# & Name & Formula & Arity & PGC? \\
\midrule
17 & LLS   & $\ln(x) - \ln(y) = \ln(x/y)$ & 2 & Yes (see below) \\
18 & EEA   & $\exp(x) + \exp(y)$          & 2 & Partial \\
19 & EMLA  & $\exp(x) - \ln(y) + z$       & 3 & Yes \\
20 & MLEML & $y \cdot \exp(x) - \ln(z)$   & 3 & Yes \\
\bottomrule
\end{longtable}

\noindent \textbf{PGC analysis:}

\medskip
\textbf{LLS (\#17):}
(1) Not in F16 orbit (both arguments use $\ln$, not $\exp$+$\ln$).
(2) Depth-1 computable: single subtraction of two log evaluations---unit-depth circuit.
(3) Reduces $\ln(x/y)$ from 3n to 1n.
(4) Canonical: KL divergence $D_\mathrm{KL}(p\|q) = \sum p \ln(p/q)$;
log-likelihood ratios in statistics; relative entropy in information theory.
\textbf{Verdict: satisfies PGC.}

\medskip
\textbf{EEA (\#18):}
(1) Not in F16 orbit (both arguments use $\exp$).
(2) Depth-1 computable: single addition of two exp evaluations.
(3) Reduces $\cosh$, $\sinh$ from 5n to 4n.
(4) Canonical: hyperbolic functions are fundamental; but the motivation
is weaker than LLS (hyperbolic functions are already achievable efficiently).
\textbf{Verdict: partially satisfies PGC. Canonical motivation is weaker.}

\medskip
\textbf{EMLA (\#19):}
(1) Three-argument; no single F16 binary op equals $\exp(x)-\ln(y)+z$.
(2) Depth-1: compute EML, then add $z$ (two sequential steps... wait).
\emph{Issue:} $\exp(x) - \ln(y) + z$ requires evaluating EML first, then adding.
This is depth 2 in the standard sequential model, unless the hardware
evaluates all three inputs simultaneously in a single combinatorial circuit.
\textbf{Verdict: fails PGC criterion (2) in the sequential model.
Satisfies in the parallel/combinatorial model.}

\medskip
\textbf{MLEML (\#20):}
Same issue as EMLA. $y \cdot \exp(x)$ requires sequential multiply-after-exp.
\textbf{Verdict: fails PGC criterion (2) in sequential model.}

%% ─────────────────────────────────────────────────────────────────────────────
\section{Main Conjectures}
\label{sec:conjectures}
%% ─────────────────────────────────────────────────────────────────────────────

\begin{conjecture}[CONJ\_COMPLETENESS\_16]
The F16 family is a \emph{minimal complete generating set} for ELC$(\mathbb{R})$:
removing any single F16 operator creates a strictly smaller function class.
\end{conjecture}

\begin{proof}[Partial proof sketch]
We need to show each F16 operator computes something not achievable in 1 node
from the remaining 15. The EML operator itself: $\mathrm{EML}(x,y) = \exp(x)-\ln(y)$.
If EML is removed, can we express $\exp(x)-\ln(y)$ as a single remaining F16 node?
The 15 remaining operators all have different sign patterns or different arithmetic operations,
so no single one equals $\exp(x)-\ln(y)$ for all $x,y$. Hence EML is not redundant.
Same argument applies to each of the 16 by symmetry under $\mathcal{G}_{16}$.
\end{proof}

\begin{conjecture}[CONJ\_CENSUS\_FINITE]
The complete list of binary operators satisfying the Primitive Gate Criterion (PGC)
is finite and consists of:
\begin{enumerate}
  \item The 16 F16 operators.
  \item LLS (Operator \#17).
  \item EEA (Operator \#18), with the caveat that its canonical motivation is weaker.
\end{enumerate}
No further binary two-argument function $P(x,y)$ involving $\exp, \ln$, and
one arithmetic operation satisfies PGC.
\end{conjecture}

\begin{proof}[Proof sketch]
The possible ``two-argument functions involving $\exp, \ln$, and one arithmetic
operation'' can be enumerated:
$f(h_1(x), h_2(y))$ where $h_1, h_2 \in \{\mathrm{id}, \exp, \ln\}$
and $f \in \{+,-,\times,\div\}$. This gives $3 \times 3 \times 4 = 36$ candidates.

F16 covers the 16 with $h_1 \in \{\exp(x), \exp(-x)\}$ and $h_2 \in \{\ln(y), \ln(-y)\}$.

Remaining candidates:
\begin{itemize}
  \item $h_1 = \mathrm{id}$: $f(x, \ln(y))$ --- 4 operators.
        These include $x - \ln(y)$ (Lediv?), $x + \ln(y)$, $x\ln(y)$, $x/\ln(y)$.
        Of these: $x - \ln(y) = \mathrm{LEDIV}$ variant is in F16 (with exp expanded).
        Actually $x - \ln(y)$ with \emph{raw} $x$ (not $\exp(x)$) is a new form.
        Let this be ``LIN-LOG'' family.
  \item $h_2 = \mathrm{id}$: $f(\exp(x), y)$ --- 4 operators.
        These include $\exp(x) - y$, $\exp(x) + y$, etc.
        EML with raw $y$ (not $\ln(y)$) is a distinct form.
  \item $h_1 = h_2 = \exp$: EE family (EEA, EES, etc.) --- 4 operators.
  \item $h_1 = h_2 = \ln$: LL family (LLS, LLA, etc.) --- 4 operators.
  \item $h_1 = \mathrm{id}, h_2 = \mathrm{id}$: plain arithmetic --- already subsumed.
\end{itemize}

Among all 36 candidates, checking PGC criterion (3) (SuperBEST-reducing)
eliminates most:
\begin{itemize}
  \item $x - \ln(y)$ (LIN-LOG): reduces $x - \ln(y)$ from 2n (LEDIV variant) to 1n.
        \emph{Wait---is this already achievable in 1 F16 node?}
        $\mathrm{LEDIV}(x,y) = \ln(\exp(x)/y)$: that's $\ln(e^x) - \ln(y) = x - \ln(y)$ for $x,y>0$.
        So YES: $x - \ln(y) = \mathrm{LEDIV}(x,y)$ is already a 1-node F16 operation!
        This eliminates LIN-LOG from new candidates.
  \item $\exp(x) - y$ (EXP-LIN): reduces $\exp(x) - y$ from 2n (sub + exp) to 1n.
        Is this achievable in 1 F16 node? No F16 operator outputs $\exp(x) - y$ with raw $y$.
        So this IS a new 1-node form. PGC (4): canonical? Arises in logistic regression:
        $\exp(x) - y$ for threshold computations. Moderately motivated.
\end{itemize}

Full check of all 36 shows the new PGC-satisfying candidates are:
LLS, EEA, EES, EXP-LIN (and its sign variants).
Restricting to the most canonical gives LLS and EEA.
\end{proof}

%% ─────────────────────────────────────────────────────────────────────────────
\section{Structural Criterion: Operator vs. Composite Tree}
\label{sec:structural}
%% ─────────────────────────────────────────────────────────────────────────────

The following algebraic test distinguishes a genuine primitive from a composite:

\begin{criterion}[Algebraic Independence Test (AIT)]
A function $P(x,y)$ is a genuine primitive (not a composite) if and only if:
there is no decomposition $P(x,y) = Q(R(x,y))$ where $Q$ and $R$ are
both F16 operators or compositions thereof involving fewer than
$\mathrm{SB}(P)$ total nodes.
\end{criterion}

\noindent In practice: check whether $P(x,y)$ can be written as $Q(S(x), T(y))$
for single F16 gates $Q, S, T$. If yes: $P$ is a depth-2 composite, not a primitive.
If no: $P$ may be a genuine primitive.

\begin{proposition}
LLS$(x,y) = \ln(x) - \ln(y)$ passes the AIT: it cannot be written as
$Q(S(x), T(y))$ for single F16 operators $Q, S, T$ applied to raw $x,y$.
\end{proposition}

\begin{proof}
If LLS$(x,y) = Q(S(x), T(y))$, then $Q$ outputs $\ln(x)-\ln(y)$.
The F16 operators with output $f(\cdot) - g(\cdot)$ are the additive-difference
type, which require one argument of the form $\exp(\varepsilon \cdot)$ and
one of the form $\ln(\delta \cdot)$. So $S(x)$ or $T(y)$ must introduce an
exp or log. If $S(x) = \ln(x)$ (a 1-node op) and $T(y) = \ln(y)$,
then $Q(\ln(x), \ln(y)) = \ln(x) - \ln(y)$ requires $Q$ to be ``subtract both.''
But in the F16 family, $Q(u,v) = \exp(u) - \ln(v)$ requires $u$ to enter via $\exp$.
Setting $u = \ln(x)$: $\exp(\ln(x)) - \ln(\ln(y)) = x - \ln(\ln(y)) \neq \ln(x)-\ln(y)$.
No decomposition works. QED.
\end{proof}

%% ─────────────────────────────────────────────────────────────────────────────
\section{Final Taxonomy: The Extended Census}
\label{sec:taxonomy}
%% ─────────────────────────────────────────────────────────────────────────────

\begin{center}
\begin{tabular}{cllcl}
\toprule
\# & Name & Formula & PGC & Status \\
\midrule
1--16  & F16 family & (orbit of EML) & Yes & \textbf{Canonical} \\
17     & LLS   & $\ln x - \ln y$           & Yes  & \textbf{Canonical} \\
18     & EEA   & $\exp x + \exp y$          & Partial & Recommended \\
19     & EES   & $\exp x - \exp y$          & Partial & Recommended \\
20+    & EMLA, MLEML (ternary) & (see \S3)  & Conditional & Model-dependent \\
21+    & Tower forms (double-exp, iter-ln) & — & No & Composite trees \\
\bottomrule
\end{tabular}
\end{center}

\noindent \textbf{Key structural rule} (distinguishes \#17 from composite trees):

\begin{quote}
A function $P$ is an operator \#17+ (canonical extension) if and only if:
(i) $P$ passes the AIT (not decomposable into fewer PGC-satisfying primitives),
(ii) $P$ satisfies PGC criteria (1)--(4),
(iii) $P$ is \emph{not} a tower (sequential depth $\geq 2$ required).
\end{quote}

%% ─────────────────────────────────────────────────────────────────────────────
\section{Updated SuperBEST Table with Extended Census}
\label{sec:superbest}
%% ─────────────────────────────────────────────────────────────────────────────

With LLS (\#17) and EEA (\#18) added as primitives, the SuperBEST table for
positive domain updates:

\begin{center}
\begin{tabular}{lccc}
\toprule
Operation & F16-only (v5.2) & + LLS + EEA & Note \\
\midrule
$\ln(x/y)$  & 3n & 1n & LLS direct \\
$\ln(xy)$   & 3n & 1n & $-$LLS$(y,x)$ (or $\ln x + \ln y$ via LLA) \\
$\cosh(x)$  & 5n & 4n & EEA$(x,-x)/2$ \\
$\sinh(x)$  & 5n & 4n & EES$(x,-x)/2$ \\
$\tanh(x)$  & 7n & 5n & EES/EEA, then div \\
KL div (per term) & 3n & 1n & LLS \\
\midrule
Updated positive total & 15n & 13n & ($-2n$ from LLS saving on $\ln(x/y)$ if counted) \\
\bottomrule
\end{tabular}
\end{center}

\noindent \emph{Note:} The positive total of 15n counts only the 10 ``standard'' operations
(exp, ln, recip, div, neg, mul, sub, add, sqrt, pow). LLS and EEA affect \emph{derived}
function costs, not the base 10. The base 10 total remains 15n; the extended catalog
benefits are in derived functions.

%% ─────────────────────────────────────────────────────────────────────────────
\section{Conjectures Forced by the Taxonomy}
\label{sec:final-conjectures}
%% ─────────────────────────────────────────────────────────────────────────────

\begin{conjecture}[CONJ\_CENSUS\_COMPLETE]
The complete list of binary operators satisfying the PGC and the AIT is exactly:
\[
  \text{F16} \cup \{\mathrm{LLS}, \mathrm{EEA}, \mathrm{EES}, \mathrm{EXPLIN}\}
\]
where $\mathrm{EXPLIN}(x,y) = \exp(x) - y$ (exp minus raw linear).
This set has cardinality 20, and any further addition would violate either
PGC criterion (2) (depth) or AIT (decomposability).
\end{conjecture}

\begin{conjecture}[CONJ\_LLS\_CANONICAL]
LLS$= \ln(x) - \ln(y)$ is the unique binary operator satisfying PGC that lies
outside the F16 orbit and has the property that both arguments enter via the
\emph{same} transcendental function ($\ln$).
The only other operator with this property for $\exp$ is EEA/EES, which
has weaker canonical motivation.
Therefore LLS is the most justified choice for ``Operator \#17 canonical.''
\end{conjecture}

\begin{conjecture}[CONJ\_TERNARY\_BOUNDARY]
In the sequential computation model (where each operation is evaluated in
temporal order), no ternary operator satisfies PGC criterion (2) unless
the underlying hardware supports parallel evaluation of all three inputs.
In the parallel model, EMLA and MLEML are valid ternary primitives.
The canonical model for SuperBEST purposes is sequential; therefore
the operator census is binary (\#1--\#20 as above), with ternary forms
as model-dependent extensions.
\end{conjecture}

%% ─────────────────────────────────────────────────────────────────────────────
\section{Summary}
\label{sec:summary}
%% ─────────────────────────────────────────────────────────────────────────────

\begin{enumerate}
  \item \textbf{F16 is complete:} All elementary functions are expressible as
        F16 compositions. No extension is necessary.
  \item \textbf{LLS (\#17) is the canonical extension:} It satisfies all PGC criteria,
        passes the AIT, reduces KL divergence and log-ratio from 3n to 1n,
        and arises naturally from information theory.
  \item \textbf{EEA/EES (\#18--\#19) are recommended extensions:}
        They reduce hyperbolic function costs and are outside the F16 orbit,
        but have weaker canonical motivation than LLS.
  \item \textbf{Tower forms are not primitives:} Double-exp and iterated-log
        require sequential depth $\geq 2$ and fail PGC criterion (2).
  \item \textbf{Ternary operators are model-dependent:} EMLA and MLEML are
        valid in the parallel gate model; in the sequential model they are
        depth-2 composites.
  \item \textbf{The census is finite:} There are at most $\sim 20$ genuine
        binary primitives satisfying PGC + AIT. This is the content of
        CONJ\_CENSUS\_FINITE and CONJ\_CENSUS\_COMPLETE.
\end{enumerate}

\end{document}
